\chapter{Cha mẹ và gia đình}
\markboth{Cha mẹ và gia đình}{Parents and Family}

\section{Mẹ gọi, người con nhắn sẽ gọi lại sau.}
\begin{truyen}{Mẹ gọi, người con nhắn sẽ gọi lại sau.}{One missed call can become forever.}
\chuhoa{N}{gày nào mẹ cũng gọi chỉ để hỏi ăn cơm chưa. Người con bận, cúp máy nhanh và nghĩ tối sẽ gọi lại. Có những cuộc gọi tưởng rất bình thường, nhưng đến lúc không còn nữa, người ta mới hiểu đó là một phần bình yên lớn nhất đời mình.}
\textit{(A mother called every day just to ask whether her child had eaten. The child was busy and thought a return call could wait. Some calls seem ordinary until they disappear forever, and only then do we understand how much peace they held.)}
\end{truyen}
\begin{baihoc}
Cha mẹ già đi không ồn ào; họ lặng lẽ đi qua thời gian cho đến một ngày ta giật mình vì đã quá muộn để ở bên như trước.
\textit{(Parents grow old quietly; they move through time silently until one day we realize it is too late to be with them as before.)}
\end{baihoc}
\begin{ghinhoanh}
\textbf{Short English to remember:}

One missed call can become forever.

Call back while you still can.
\end{ghinhoanh}
\begin{tuhocnhanh}
1. Lâu rồi em chưa ngồi nói chuyện thật sự với cha mẹ về điều gì? \textit{(What have you not truly talked about with your parents for a long time?)}

2. Nếu tối nay gọi về nhà, em muốn nói câu nào trước tiên? \textit{(If you called home tonight, what would you want to say first?)}
\end{tuhocnhanh}
\ngancach

\section{Người con chê món ăn mẹ nấu vì nghĩ hôm nào cũng có.}
\begin{truyen}{Người con chê món ăn mẹ nấu vì nghĩ hôm nào cũng có.}{Ordinary love is easy to overlook.}
\chuhoa{T}{rên bàn cơm, cậu chỉ thấy món quen và những lời nhắc quen. Ra ở riêng rồi, cậu mới hiểu bữa cơm lặp lại ấy chính là tình thương ổn định nhất mà đời mình từng có.}
\textit{(At the dinner table, he saw only familiar dishes and familiar reminders. After moving out, he realized that those repetitive meals were the most stable love he had ever known.)}
\end{truyen}
\begin{baihoc}
Điều ta hay chê là điều ta tưởng sẽ luôn còn; đến lúc mất đi mới biết nó quý vì đã từng quá quen.
\textit{(What we often dismiss is what we assume will always remain; once it is gone, we see its worth precisely because it was so familiar.)}
\end{baihoc}
\begin{ghinhoanh}
\textbf{Short English to remember:}

Ordinary love is easy to overlook.

Familiar does not mean cheap.
\end{ghinhoanh}
\begin{tuhocnhanh}
1. Có điều gì trong gia đình em đã quen tới mức em quên biết ơn? \textit{(What in your family has become so familiar that you forget to appreciate it?)}

2. Em có thể thể hiện sự trân trọng đó bằng việc nhỏ nào hôm nay? \textit{(How can you show appreciation for it in a small way today?)}
\end{tuhocnhanh}
\ngancach

\section{Người trẻ nghĩ cha mẹ sẽ hiểu mình mà không cần nói.}
\begin{truyen}{Người trẻ nghĩ cha mẹ sẽ hiểu mình mà không cần nói.}{Silence also creates distance.}
\chuhoa{C}{ô giữ hết áp lực trong lòng, chỉ đáp cụt lủn với gia đình. Khi hiểu lầm chồng lên hiểu lầm, cô mới biết tình thân cũng cần được giải thích và chăm sóc, không tự bền mãi nếu ai cũng im lặng.}
\textit{(She kept all her pressure inside and answered her family in short, cold replies. When misunderstandings piled up, she learned that family bonds also need explanation and care; they do not stay strong by themselves if everyone is silent.)}
\end{truyen}
\begin{baihoc}
Yêu nhau mà không nói ra điều cần nói cũng có thể đẩy nhau xa dần.
\textit{(Even people who love each other can drift apart if they never say what needs to be said.)}
\end{baihoc}
\begin{ghinhoanh}
\textbf{Short English to remember:}

Silence also creates distance.

Love needs honest words too.
\end{ghinhoanh}
\begin{tuhocnhanh}
1. Có hiểu lầm nào trong gia đình em vẫn còn bỏ đó? \textit{(Is there any family misunderstanding still left unresolved?)}

2. Em sợ điều gì nhất khi mở lời thật lòng với người thân? \textit{(What do you fear most about speaking honestly with family?)}
\end{tuhocnhanh}
\ngancach

\section{Một người chỉ đến bệnh viện khi cha yếu hẳn.}
\begin{truyen}{Một người chỉ đến bệnh viện khi cha yếu hẳn.}{Visit before the final urgency.}
\chuhoa{N}{hững lần cha nằm viện nhẹ, anh viện cớ bận. Đến lúc bệnh nặng, anh có mặt đủ cả ngày nhưng trong lòng không còn yên vì biết mình đã vắng ở những ngày cha còn cần một sự quan tâm bình thường.}
\textit{(When his father was mildly ill, he said he was too busy to visit. When the illness became serious, he stayed all day, but could not find peace because he knew he had been absent on the ordinary days that still mattered.)}
\end{truyen}
\begin{baihoc}
Tận tâm phút cuối không xóa được sự thờ ơ kéo dài từ trước đó.
\textit{(Deep care at the final moment does not erase long seasons of earlier neglect.)}
\end{baihoc}
\begin{ghinhoanh}
\textbf{Short English to remember:}

Visit before the final urgency.

Care should not wait for crisis.
\end{ghinhoanh}
\begin{tuhocnhanh}
1. Em có đang chỉ quan tâm người thân khi có chuyện lớn xảy ra không? \textit{(Do you only pay attention to loved ones when something serious happens?)}

2. Một sự hiện diện bình thường mà em có thể làm tuần này là gì? \textit{(What ordinary act of presence can you offer this week?)}
\end{tuhocnhanh}
\ngancach

\section{Người con thành đạt nhưng không còn ngồi ăn cơm với nhà.}
\begin{truyen}{Người con thành đạt nhưng không còn ngồi ăn cơm với nhà.}{Success cannot replace belonging.}
\chuhoa{C}{ậu có tiền, có xe, có lịch làm việc dày đặc. Nhưng mỗi lần về nhà chỉ ghé chốc lát. Đến khi nhìn bàn ăn thưa dần người cũ, cậu mới hiểu thành công ngoài xã hội không tự động bù được sự vắng mặt trong gia đình.}
\textit{(He had money, a car, and a packed schedule. But each visit home was brief. When he saw fewer familiar faces around the table, he understood that success outside cannot automatically make up for absence inside the family.)}
\end{truyen}
\begin{baihoc}
Một đời đi xa mà không còn chỗ để quay về cũng là một kiểu nghèo.
\textit{(A life that goes far but no longer has a real place to return to is a kind of poverty.)}
\end{baihoc}
\begin{ghinhoanh}
\textbf{Short English to remember:}

Success cannot replace belonging.

Come home before home changes.
\end{ghinhoanh}
\begin{tuhocnhanh}
1. Em đang đầu tư nhiều hơn cho thành tựu hay cho sự gắn bó gia đình? \textit{(Are you investing more in achievement or in family connection?)}

2. Điều gì làm em thấy thật sự thuộc về gia đình mình? \textit{(What makes you truly feel you belong in your family?)}
\end{tuhocnhanh}
\ngancach